\documentclass[11pt,a4paper]{article}
\usepackage[utf8]{inputenc}
\usepackage[T1]{fontenc}
\usepackage[french]{babel}
\usepackage{lmodern}
\usepackage{amsmath,amssymb}
\usepackage[margin=2.3cm]{geometry}
\usepackage{microtype}
\usepackage{hyperref}
\usepackage{booktabs}
\usepackage{graphicx}
\usepackage{pgfplots}
\pgfplotsset{compat=1.18}
\usepackage{listings}
\lstset{basicstyle=\ttfamily\small, frame=single, breaklines=true,
  columns=fullflexible, keepspaces=true}
\title{Projet MLSF : résolution du jeu Meowdoku\\
  \large Question 5 : analyse des performances des heuristiques}
\author{}
\date{}

% Courbe « instances résolues en fonction du temps CPU » d'une heuristique.
\newcommand{\cactus}[3]{\addplot[mark=none, thick, #3] table[x=cpu, y=solved]
  {generated/cactus-#1-#2.dat}; \addlegendentry{#2}}

\begin{document}
\maketitle

\section{Objectif et protocole}

La question demande d'instrumenter les solveurs DPLL$_H$ de la question 4,
de les comparer sur des instances variées (forme, statut SAT/UNSAT,
Meowdoku), puis de répondre à trois questions : les résultats sont-ils
conformes à nos intuitions, combien une heuristique améliore-t-elle
l'algorithme 2, et le recalcul dynamique vaut-il son coût ?

Toutes les mesures sont produites par \texttt{bench\_q5.py}, en deux étapes :
\begin{lstlisting}
uv run python bench_q5.py run --timeout 60    # environ 10 minutes, 6 processus
uv run python bench_q5.py report              # tableaux et courbes de ce rapport
\end{lstlisting}
La première étape écrit une ligne par résolution dans
\texttt{results/q5\_runs.csv} ; la seconde produit les fichiers du dossier
\texttt{rapports/question5/generated}, inclus directement dans ce rapport.
Les tableaux ne sont donc jamais recopiés à la main.

\subsection{Instances}

Nous avons retenu 173 instances réparties en cinq familles, choisies pour
couvrir les trois axes demandés :
\begin{itemize}
  \item \textbf{3-SAT aléatoire SAT et UNSAT} : 30 instances uf50 et 20 uf75
    (SAT), 30 uuf50 et 20 uuf75 (UNSAT) de SATLIB, au seuil de
    satisfiabilité ($m/n \approx 4{,}26$), là où le 3-SAT aléatoire est le
    plus difficile. Elles n'ont aucune structure : c'est le cas où une
    heuristique ne peut s'appuyer que sur des statistiques d'occurrences.
  \item \textbf{Instances structurées} : les 23 instances publiques des
    enseignants qui ne sont pas du 3-SAT aléatoire à 50 variables
    (coloration de graphes, planification, circuits, pigeons,
    quasigroupes\dots), de 42 à 1363 variables ; 17 SAT et 6 UNSAT.
  \item \textbf{Meowdoku (question 1)} : 20 grilles produites par notre
    générateur, de taille $n = 8$, 10, 12, 14 et 16 (4 par taille), SAT
    par construction.
  \item \textbf{Meowdoku à régions aléatoires} : 30 grilles de taille 6 à
    10 dont les $n$ régions connexes sont grandies depuis $n$ cases tirées
    au hasard, sans solution placée à l'avance (fonction
    \texttt{random\_regions\_grid}).
    Elles fournissent des grilles Meowdoku \textbf{UNSAT}, que le générateur
    de la question 1 ne produit pas : 7 sont SAT et 23 UNSAT.
\end{itemize}
Le statut de chaque instance est connu (SATLIB, \texttt{expected.json},
construction) sauf pour les grilles à régions aléatoires, dont le statut est
celui renvoyé par les heuristiques. Le script s'arrête si deux heuristiques
se contredisent ou si un modèle renvoyé ne satisfait pas la formule :
aucune contradiction ni aucun modèle faux sur les 2768 résolutions.

\subsection{Critères mesurés}

Les quatre compteurs de l'énoncé sont ceux de la classe \texttt{Stats}
de \texttt{dpll.py}, incrémentés aux lignes correspondantes de
l'algorithme : décisions (ligne 17), retours arrière (ligne 21),
propagations unitaires (ligne 5) et littéraux purs (ligne 11). S'y ajoutent :
\begin{itemize}
  \item le \textbf{temps CPU} (\texttt{time.process\_time}) ;
  \item la \textbf{part du temps passée dans l'heuristique} : temps cumulé
    des appels à \texttt{choose}, divisé par la durée de la résolution ;
  \item trois critères de notre choix :
    \begin{itemize}
      \item le \textbf{nombre d'instances résolues} dans le délai de 60 s,
        celui de la compétition ;
      \item le \textbf{score PAR-2}, temps moyen où une instance non résolue
        compte pour deux fois le délai (120 s). C'est la mesure des
        compétitions SAT : elle combine vitesse et robustesse ;
      \item le nombre de \textbf{propagations unitaires par décision},
        qui mesure ce qu'une décision \og rapporte \fg{}.
    \end{itemize}
\end{itemize}

\subsection{Choix méthodologiques}

\paragraph{Compteurs sur les instances résolues par toutes les
heuristiques.} Une résolution interrompue a des compteurs partiels :
comparer ses 10\,000 décisions aux 50 d'une heuristique qui a conclu n'a pas
de sens. Les médianes des compteurs sont donc calculées sur les instances
résolues par les 16 heuristiques du groupe ; le nombre de résolues et le
PAR-2 portent, eux, sur toutes les instances.

\paragraph{Médianes et moyennes géométriques.} Les temps et les compteurs
varient sur plusieurs ordres de grandeur d'une instance à l'autre : une
moyenne arithmétique serait dominée par une ou deux instances. Nous
utilisons des médianes et, pour comparer une heuristique à H0, la
\textbf{moyenne géométrique des rapports} $\text{H0}/H$ calculés instance par
instance : c'est la moyenne adaptée à des rapports (un gain $\times 2$ et une
perte $\times 2$ se compensent). Pour les décisions on utilise
$(d_{\text{H0}}+1)/(d_H+1)$, car certaines instances n'en demandent aucune.

\paragraph{Une graine par instance pour H0.} H0 est aléatoire. Sa graine
dérive du nom de l'instance : chaque instance a son propre tirage et la
campagne est reproductible. La variabilité due à la graine a été étudiée à
la question 2 ; elle est ici moyennée sur les nombreuses instances de
chaque groupe.

\paragraph{Parallélisme.} Les 2768 résolutions sont réparties sur 6
processus (\texttt{multiprocessing}). Le temps CPU est mesuré dans chaque
processus et ne compte donc pas le travail des autres. Les plus gros
fichiers passent en premier pour éviter qu'une longue résolution lancée
en dernier retarde la fin de la campagne.

\section{Résultats par groupe}

Les familles qui mélangent les deux statuts (instances structurées,
grilles aléatoires) sont séparées en SAT et UNSAT, comme le suggère
l'énoncé :
\begin{center}
\begin{tabular}{lrr}
\toprule
Groupe & instances & résolues par toutes \\
\midrule
3-SAT aléatoire SAT & 50 & 50 \\
3-SAT aléatoire UNSAT & 50 & 50 \\
Instances structurées (SAT) & 17 & 12 \\
Instances structurées (UNSAT) & 6 & 5 \\
Meowdoku (question 1) & 20 & 19 \\
Meowdoku, régions aléatoires (SAT) & 7 & 7 \\
Meowdoku, régions aléatoires (UNSAT) & 23 & 23 \\
\bottomrule
\end{tabular}

\end{center}

Pour chaque groupe, le tableau donne le nombre d'instances résolues, le
PAR-2, puis les médianes des compteurs (sur les instances résolues par
toutes), les propagations unitaires par décision et la part du temps
passée dans l'heuristique.

\newcommand{\detail}[2]{%
  \begin{table}[htbp]\centering\small
  \input{generated/detail-#1.tex}
  \caption{#2}\end{table}}

\detail{uf}{3-SAT aléatoire SAT (uf50, uf75).}
\detail{uuf}{3-SAT aléatoire UNSAT (uuf50, uuf75).}
\detail{struct-sat}{Instances structurées SAT.}
\detail{struct-unsat}{Instances structurées UNSAT. Le PAR-2 est dominé par
  public-03, qu'aucune heuristique ne résout.}
\detail{meow}{Grilles Meowdoku de la question 1 (SAT).}
\detail{meowr-sat}{Grilles Meowdoku à régions aléatoires, SAT.}
\detail{meowr-unsat}{Grilles Meowdoku à régions aléatoires, UNSAT.}

\clearpage
\section{Gain par rapport à l'algorithme 2}

Les deux tableaux suivants résument les précédents : pour chaque groupe, la
moyenne géométrique du rapport $\text{H0}/H$ sur les instances résolues par
toutes les heuristiques. Une valeur supérieure à 1 signifie que $H$ fait
mieux que l'algorithme 2. Entre parenthèses, le nombre d'instances
résolues par $H$ en moins de 60 s.

\begin{table}[htbp]\centering\footnotesize\setlength{\tabcolsep}{4pt}
\begin{tabular}{lrrrrrrr}
\toprule
H & \rotatebox{70}{3-SAT aléatoire SAT} & \rotatebox{70}{3-SAT aléatoire UNSAT} & \rotatebox{70}{Instances structurées (SAT)} & \rotatebox{70}{Instances structurées (UNSAT)} & \rotatebox{70}{Meowdoku (question 1)} & \rotatebox{70}{Meowdoku, régions aléatoires (SAT)} & \rotatebox{70}{Meowdoku, régions aléatoires (UNSAT)} \\
\midrule
\texttt{H0} & 1{,}0 (50) & 1{,}0 (50) & 1{,}0 (15) & 1{,}0 (5) & 1{,}0 (20) & 1{,}0 (7) & 1{,}0 (23) \\
\texttt{H1} & 1{,}0 (50) & 1{,}1 (50) & 1{,}2 (16) & 3{,}2 (5) & 2{,}2 (19) & 1{,}3 (7) & 1{,}3 (23) \\
\texttt{H2s} & 6{,}8 (50) & 7{,}9 (50) & 2{,}4 (14) & 3{,}8 (5) & 0{,}9 (20) & 0{,}6 (7) & 0{,}5 (23) \\
\texttt{H2d} & 5{,}6 (50) & 6{,}8 (50) & 1{,}7 (15) & 4{,}6 (5) & 0{,}8 (20) & 0{,}4 (7) & 0{,}7 (23) \\
\texttt{H3s} & 7{,}7 (50) & 5{,}6 (50) & 2{,}4 (16) & 3{,}0 (5) & 0{,}8 (20) & 0{,}7 (7) & 0{,}5 (23) \\
\texttt{H3d} & 6{,}7 (50) & 4{,}8 (50) & 1{,}4 (16) & 2{,}6 (5) & 0{,}8 (20) & 0{,}3 (7) & 0{,}7 (23) \\
\texttt{H4s} & 7{,}2 (50) & 7{,}9 (50) & 2{,}9 (16) & 25{,}5 (5) & 0{,}8 (20) & 0{,}8 (7) & 0{,}7 (23) \\
\texttt{H4d} & 11{,}0 (50) & 14{,}6 (50) & 3{,}3 (17) & 23{,}6 (5) & 0{,}8 (20) & 0{,}3 (7) & 0{,}7 (23) \\
\texttt{H5s} & 8{,}0 (50) & 5{,}9 (50) & 3{,}2 (16) & 22{,}6 (5) & 0{,}8 (20) & 0{,}7 (7) & 0{,}7 (23) \\
\texttt{H5d} & 8{,}5 (50) & 9{,}5 (50) & 1{,}8 (17) & 15{,}7 (5) & 0{,}8 (20) & 0{,}3 (7) & 0{,}7 (23) \\
\texttt{H6s} & 7{,}1 (50) & 8{,}0 (50) & 2{,}8 (16) & 22{,}3 (5) & 0{,}8 (20) & 0{,}5 (7) & 0{,}5 (23) \\
\texttt{H6d} & 8{,}8 (50) & 10{,}9 (50) & 2{,}2 (16) & 18{,}6 (5) & 0{,}7 (20) & 0{,}4 (7) & 0{,}7 (23) \\
\texttt{H7s} & 7{,}8 (50) & 5{,}7 (50) & 2{,}5 (15) & 15{,}5 (5) & 0{,}8 (20) & 0{,}7 (7) & 0{,}5 (23) \\
\texttt{H7d} & 8{,}0 (50) & 7{,}3 (50) & 1{,}3 (16) & 13{,}3 (5) & 0{,}7 (20) & 0{,}3 (7) & 0{,}7 (23) \\
\texttt{H8s} & 8{,}2 (50) & 5{,}2 (50) & 3{,}7 (15) & 20{,}6 (5) & 3{,}3 (20) & 1{,}2 (7) & 1{,}7 (23) \\
\texttt{H8d} & 7{,}2 (50) & 6{,}3 (50) & 2{,}9 (17) & 14{,}9 (5) & 2{,}7 (20) & 1{,}0 (7) & 2{,}1 (23) \\
\bottomrule
\end{tabular}

\caption{Gain en temps CPU par rapport à H0 (nombre d'instances résolues).}
\label{tab:gain-temps}
\end{table}
\begin{table}[htbp]\centering\footnotesize\setlength{\tabcolsep}{4pt}
\begin{tabular}{lrrrrrrr}
\toprule
H & \rotatebox{70}{3-SAT aléatoire SAT} & \rotatebox{70}{3-SAT aléatoire UNSAT} & \rotatebox{70}{Instances structurées (SAT)} & \rotatebox{70}{Instances structurées (UNSAT)} & \rotatebox{70}{Meowdoku (question 1)} & \rotatebox{70}{Meowdoku, régions aléatoires (SAT)} & \rotatebox{70}{Meowdoku, régions aléatoires (UNSAT)} \\
\midrule
\texttt{H0} & 1{,}0 & 1{,}0 & 1{,}0 & 1{,}0 & 1{,}0 & 1{,}0 & 1{,}0 \\
\texttt{H1} & 1{,}0 & 1{,}1 & 0{,}9 & 3{,}5 & 5{,}6 & 1{,}2 & 1{,}8 \\
\texttt{H2s} & 6{,}6 & 9{,}2 & 2{,}7 & 4{,}8 & 0{,}6 & 0{,}4 & 0{,}5 \\
\texttt{H2d} & 8{,}0 & 11{,}5 & 2{,}0 & 8{,}2 & 1{,}2 & 0{,}3 & 0{,}6 \\
\texttt{H3s} & 6{,}1 & 6{,}2 & 2{,}7 & 3{,}4 & 0{,}6 & 0{,}4 & 0{,}5 \\
\texttt{H3d} & 7{,}5 & 7{,}3 & 1{,}4 & 3{,}8 & 1{,}2 & 0{,}3 & 0{,}6 \\
\texttt{H4s} & 6{,}6 & 9{,}2 & 2{,}4 & 30{,}3 & 0{,}6 & 0{,}4 & 0{,}8 \\
\texttt{H4d} & 12{,}9 & 22{,}0 & 2{,}7 & 36{,}8 & 1{,}3 & 0{,}4 & 0{,}8 \\
\texttt{H5s} & 6{,}1 & 6{,}2 & 2{,}3 & 30{,}8 & 0{,}6 & 0{,}4 & 0{,}7 \\
\texttt{H5d} & 8{,}4 & 13{,}0 & 1{,}3 & 28{,}1 & 1{,}2 & 0{,}3 & 0{,}6 \\
\texttt{H6s} & 6{,}6 & 9{,}2 & 3{,}1 & 33{,}1 & 0{,}5 & 0{,}3 & 0{,}5 \\
\texttt{H6d} & 12{,}7 & 19{,}7 & 2{,}1 & 39{,}7 & 1{,}2 & 0{,}4 & 0{,}7 \\
\texttt{H7s} & 6{,}1 & 6{,}2 & 2{,}3 & 26{,}0 & 0{,}6 & 0{,}4 & 0{,}5 \\
\texttt{H7d} & 9{,}3 & 12{,}2 & 1{,}2 & 27{,}5 & 1{,}2 & 0{,}3 & 0{,}6 \\
\texttt{H8s} & 6{,}7 & 5{,}5 & 3{,}3 & 32{,}2 & 8{,}9 & 1{,}9 & 4{,}4 \\
\texttt{H8d} & 9{,}4 & 10{,}3 & 3{,}6 & 27{,}5 & 9{,}3 & 2{,}0 & 5{,}0 \\
\bottomrule
\end{tabular}

\caption{Gain en nombre de décisions par rapport à H0.}
\label{tab:gain-decisions}
\end{table}

\begin{figure}[htbp]\centering
\begin{tikzpicture}
\begin{semilogxaxis}[width=0.92\linewidth, height=7.5cm, xlabel={temps CPU (s)},
  ylabel={instances résolues}, legend pos=south east, legend columns=2,
  legend style={font=\small}, grid=major, xmin=1e-3, xmax=60]
  \cactus{satlib}{H0}{black}
  \cactus{satlib}{H1}{gray}
  \cactus{satlib}{H2d}{orange}
  \cactus{satlib}{H4s}{blue, dashed}
  \cactus{satlib}{H4d}{blue}
  \cactus{satlib}{H6d}{red}
  \cactus{satlib}{H8d}{green!60!black}
\end{semilogxaxis}
\end{tikzpicture}
\caption{Nombre d'instances résolues (3-SAT aléatoire et structurées, 123
  instances) en fonction du temps CPU accordé à chacune. Plus une courbe est
  haute et à gauche, meilleure est l'heuristique.}
\label{fig:cactus-satlib}
\end{figure}

\paragraph{3-SAT aléatoire.} C'est là que les heuristiques apportent le
plus : toutes celles fondées sur les occurrences (H2 à H8) divisent le
temps de 5 à 15 et le nombre de décisions de 5 à 22. La meilleure est
\textbf{H4d} (MOMS dynamique) : 17 décisions médianes contre 249 pour H0
sur les instances SAT, 40 contre 568 sur les UNSAT, soit un gain en temps
de 11 et 15. H6d (Jeroslow-Wang dynamique) la suit de près. H1 ne fait
pas mieux que le hasard : l'ordre des numéros de variables n'a aucun sens
sur une formule aléatoire.

\paragraph{SAT contre UNSAT.} Les gains sont plus forts sur les instances
UNSAT (jusqu'à 22 en décisions contre 13 en SAT). C'est attendu : pour
prouver l'insatisfiabilité, DPLL doit parcourir tout l'arbre de recherche,
et une bonne heuristique en réduit la taille entière. Sur une instance SAT,
même un mauvais choix peut tomber rapidement sur une solution, ce qui
réduit l'écart. Sur les UNSAT, décisions et retours arrière sont égaux :
chaque décision finit par être réfutée.

\paragraph{Instances structurées.} Les gains en temps sont plus modestes en
SAT (de 1,3 à 3,7) mais spectaculaires en UNSAT (jusqu'à 25 pour H4s) : les
heuristiques sensibles à la taille des clauses (H4 à H8) font 26 à 40 fois
moins de décisions que H0, contre 3 à 8 fois pour H2 et H3. Ce sont aussi les seules qui
résolvent toutes les instances SAT structurées : \textbf{H4d, H5d et H8d}
(17/17). La courbe de la figure~\ref{fig:cactus-satlib} montre que H4d domine
les autres sur l'ensemble des 123 instances non Meowdoku : avec 0,1 s par
instance elle en résout 111, contre 50 pour H0 ; H4d, H5d et H8d en
résolvent 122 en 60 s.
Aucune heuristique ne résout public-03 (UNSAT, 63 variables), une des
deux instances annoncées comme difficiles : réduire l'arbre ne suffit pas,
il faudra une structure de données plus efficace (question 6).

\section{Meowdoku}

\begin{table}[htbp]\centering\small
\begin{tabular}{lrrrrr}
\toprule
H & $n=8$ & $n=10$ & $n=12$ & $n=14$ & $n=16$ \\
\midrule
\texttt{H0} & 0{,}003 & 0{,}025 & 0{,}020 & 0{,}048 & 0{,}530 \\
\texttt{H1} & 0{,}005 & 0{,}013 & 0{,}022 & 0{,}040 & 0{,}603 \\
\texttt{H2s} & 0{,}005 & 0{,}018 & 0{,}069 & 0{,}181 & 0{,}204 \\
\texttt{H2d} & 0{,}013 & 0{,}033 & 0{,}089 & 0{,}168 & 0{,}328 \\
\texttt{H3s} & 0{,}005 & 0{,}020 & 0{,}075 & 0{,}185 & 0{,}203 \\
\texttt{H3d} & 0{,}015 & 0{,}037 & 0{,}095 & 0{,}173 & 0{,}321 \\
\texttt{H4s} & 0{,}005 & 0{,}014 & 0{,}069 & 0{,}191 & 0{,}716 \\
\texttt{H4d} & 0{,}011 & 0{,}037 & 0{,}092 & 0{,}183 & 0{,}313 \\
\texttt{H5s} & 0{,}005 & 0{,}016 & 0{,}070 & 0{,}183 & 0{,}706 \\
\texttt{H5d} & 0{,}012 & 0{,}036 & 0{,}088 & 0{,}166 & 0{,}310 \\
\texttt{H6s} & 0{,}005 & 0{,}017 & 0{,}097 & 0{,}607 & 0{,}217 \\
\texttt{H6d} & 0{,}012 & 0{,}041 & 0{,}094 & 0{,}200 & 0{,}372 \\
\texttt{H7s} & 0{,}005 & 0{,}015 & 0{,}068 & 0{,}365 & 0{,}267 \\
\texttt{H7d} & 0{,}012 & 0{,}037 & 0{,}098 & 0{,}190 & 0{,}356 \\
\texttt{H8s} & 0{,}003 & 0{,}006 & 0{,}017 & 0{,}029 & 0{,}068 \\
\texttt{H8d} & 0{,}004 & 0{,}010 & 0{,}019 & 0{,}043 & 0{,}073 \\
\bottomrule
\end{tabular}

\caption{Temps CPU médian (s) par taille de grille, Meowdoku de la
  question 1 (4 grilles par taille).}
\label{tab:taille}
\end{table}

\begin{figure}[htbp]\centering
\begin{tikzpicture}
\begin{semilogxaxis}[width=0.92\linewidth, height=7cm, xlabel={temps CPU (s)},
  ylabel={instances résolues}, legend pos=south east, legend columns=2,
  legend style={font=\small}, grid=major, xmin=1e-3, xmax=60]
  \cactus{meowdoku}{H0}{black}
  \cactus{meowdoku}{H1}{gray}
  \cactus{meowdoku}{H2s}{orange, dashed}
  \cactus{meowdoku}{H2d}{orange}
  \cactus{meowdoku}{H6d}{red}
  \cactus{meowdoku}{H8s}{green!60!black, dashed}
  \cactus{meowdoku}{H8d}{green!60!black}
\end{semilogxaxis}
\end{tikzpicture}
\caption{Nombre de grilles Meowdoku résolues (50 grilles) en fonction du
  temps CPU accordé à chacune.}
\label{fig:cactus-meowdoku}
\end{figure}

\paragraph{H8 est nettement la meilleure.} Sur les grilles de la question 1,
H8s et H8d sont environ 3 fois plus rapides que H0 en moyenne géométrique et
\textbf{7 à 8 fois plus rapides} sur les grilles $16 \times 16$ (0,07 s contre
0,53 s en médiane, tableau~\ref{tab:taille}). Leur temps croît beaucoup
moins vite avec la taille que celui de toutes les autres heuristiques. Une
décision de H8 (placer un chat dans la plus petite région) déclenche environ
16 propagations unitaires, contre 9 pour H0 et 0,2 pour H2d.

\paragraph{H2 à H7 : beaucoup de décisions, mais aucun retour arrière.}
La question 4 avait montré que, sur les petites grilles, H2 à H7 font plus de
décisions que le hasard (115 décisions médianes pour H2d contre 54 pour H0).
Les nouveaux compteurs nuancent ce constat : ces heuristiques ne font
\textbf{aucun} retour arrière (médiane 0 pour les versions dynamiques,
contre 43 pour H0). Elles affectent à 0 une case très contrainte après
l'autre (\og pas de chat ici \fg{}), jusqu'à ce que la clause
\og au moins un chat \fg{} d'une région devienne unitaire. Chaque décision
rapporte peu (0,2 propagation par décision) mais n'est presque jamais
remise en cause. Sur les grilles $16 \times 16$, cette stratégie prudente
rattrape H0, dont les mauvais choix aléatoires coûtent cher : elles ont un
PAR-2 bien meilleur (0,12 s contre 2,03 s), même si leur moyenne
géométrique reste légèrement inférieure à celle de H0 (0,7 à 0,9) à cause
des petites grilles. Le nombre de décisions seul est donc un indicateur
trompeur : il faut le lire avec le nombre de retours arrière.

\paragraph{Grilles UNSAT.} Les 23 grilles à régions aléatoires UNSAT sont
faciles : H0 et H8 concluent en une ou deux décisions, la contradiction
apparaissant presque immédiatement par propagation (par exemple deux
petites régions adjacentes). La plupart des variantes de H2 à H7 font au
contraire une vingtaine de décisions, toutes réfutées (sauf H4s et H5s, qui
concluent en une décision) : en choisissant \og pas de chat \fg{}, elles
retardent la propagation qui ferait apparaître la contradiction.
H8 est ici 4 à 5 fois plus économe en décisions que H0 et 2 fois plus
rapide.

\section{Statique ou dynamique ?}

Le tableau suivant compare les deux variantes de chaque heuristique sur
toutes les instances résolues par les deux. \og décisions s/d \fg{} est le
rapport (moyenne géométrique) du nombre de décisions de la variante
statique à celui de la variante dynamique : au-dessus de 1, le dynamique
choisit mieux. \og coût d'un choix d/s \fg{} compare le temps moyen d'un
appel à l'heuristique. \og temps s/d \fg{} est le rapport des temps CPU :
au-dessus de 1, le dynamique est plus rapide au total.
\begin{center}\small
\begin{tabular}{lrrrrr}
\toprule
H & résolues (s / d) & décisions s/d & coût d'un choix d/s & temps s/d & \% heur. (s / d) \\
\midrule
H2 & 169 / 170 & 1{,}28 & 6{,}3 & 0{,}87 & 7 / 20 \\
H3 & 171 / 171 & 1{,}19 & 5{,}7 & 0{,}85 & 8 / 21 \\
H4 & 171 / 172 & 1{,}74 & 5{,}0 & 1{,}31 & 8 / 40 \\
H5 & 171 / 172 & 1{,}33 & 4{,}1 & 1{,}04 & 8 / 30 \\
H6 & 171 / 171 & 1{,}79 & 7{,}0 & 1{,}17 & 8 / 30 \\
H7 & 170 / 171 & 1{,}52 & 6{,}5 & 1{,}02 & 7 / 26 \\
H8 & 170 / 172 & 1{,}34 & 6{,}4 & 0{,}97 & 7 / 32 \\
\bottomrule
\end{tabular}

\end{center}

\paragraph{Lecture.} Le recalcul dynamique améliore \emph{toujours} la
qualité des choix : la variante statique fait de 1,2 (H3) à 1,8 (H6) fois
plus de décisions que la variante dynamique. Mais
chaque choix coûte 4 à 7 fois plus cher, et la part du temps passée dans
l'heuristique passe d'environ 8\,\% à 20--40\,\%. Le bilan dépend donc du
gain en qualité :
\begin{itemize}
  \item pour \textbf{H4 et H6} (MOMS et Jeroslow-Wang), le gain en
    décisions est grand ($\times 1{,}7$ à $1{,}8$) et compense largement le
    coût : la version dynamique est 17 à 31\,\% plus rapide ;
  \item pour \textbf{H2 et H3} (fréquences brutes), le gain est faible
    ($\times 1{,}2$ à $1{,}3$) et la version dynamique est 13 à 15\,\% plus
    lente ;
  \item H5, H7 et H8 sont à l'équilibre.
\end{itemize}
C'est cohérent avec l'intuition : H4 et H6 privilégient les clauses courtes,
dont la taille change à chaque simplification ; un classement calculé sur
la formule initiale devient vite faux. Les fréquences brutes de H2 et H3
évoluent moins, et les recalculer rapporte peu. Les variantes dynamiques
résolvent aussi plus d'instances difficiles (public-12 n'est résolue par
aucune variante statique, ni par H0 ou H1, mais par toutes les
dynamiques). Sur
Meowdoku, le coût du recalcul atteint 75\,\% du temps pour H2d--H7d, et
H8s y est la plus rapide.

\paragraph{Remarque sur le coût.} Notre recalcul parcourt toute la formule
à chaque décision. Le coût d'un choix dynamique (4 à 7 fois celui d'un choix
statique) est donc une borne haute : une mise à jour incrémentale des
compteurs, faite pendant la simplification, le réduirait. C'est une piste
pour la question 6.

\section{Réponses aux questions de l'énoncé}

\paragraph{Les résultats sont-ils conformes à nos intuitions ?} En grande
partie :
\begin{itemize}
  \item les heuristiques fondées sur les occurrences battent largement le
    hasard sur le 3-SAT aléatoire, et plus encore sur les instances UNSAT,
    comme attendu ;
  \item favoriser les clauses courtes (H4, H6) est le meilleur critère
    général, conformément à l'intuition de l'énoncé (les clauses courtes
    deviennent unitaires et déclenchent la propagation) ;
  \item H1 n'apporte rien, ce qui était prévisible.
\end{itemize}
Deux résultats nous ont surpris. D'abord, sur Meowdoku, H2 à H7 choisissent
systématiquement la mauvaise polarité, ce qui nous a conduits à concevoir
H8 à la question 4. Ensuite, ce mauvais choix n'est pas catastrophique :
il donne beaucoup de décisions mais aucun retour arrière, et finit par
battre H0 sur les grandes grilles. Enfin, aucune heuristique n'est la
meilleure partout : H4d domine le 3-SAT aléatoire, H8 domine Meowdoku, et
H4d, H5d et H8d sont les seules à résoudre toutes les instances
structurées SAT.

\paragraph{De combien une heuristique améliore-t-elle l'algorithme 2 ?}
Sur le 3-SAT aléatoire, d'un facteur 5 à 15 en temps et 5 à 22 en
décisions ; sur les instances structurées UNSAT, jusqu'à 25 en temps et 40
en décisions ; sur Meowdoku, d'un facteur 3 en moyenne et 8 sur les plus
grandes grilles avec H8. Elles résolvent aussi des instances que
l'algorithme 2 ne résout pas en 60 s (public-12 pour toutes les variantes
dynamiques, public-06 pour la plupart des heuristiques). Le gain n'est pas
universel : sur les petites grilles Meowdoku, H2 à H7 sont légèrement plus
lentes que H0, et H2s ou H2d échouent sur public-04 ou public-21 que H0
résout.

\paragraph{Le recalcul dynamique vaut-il son coût ?} Il améliore toujours
la qualité des choix (1,2 à 1,8 fois moins de décisions) mais multiplie
le coût d'un choix par 4 à 7. Il est rentable pour les heuristiques qui
tiennent compte de la taille des clauses (H4, H6), dont le classement
change beaucoup au fil de la recherche ; il ne l'est pas pour les
fréquences brutes (H2, H3) ni sur Meowdoku, où le calcul de l'heuristique
représente jusqu'à 75\,\% du temps.

\paragraph{Conséquences pour la question 6.} Pour le solveur optimisé de
Meowdoku, nous retenons H8, qui est aussi parmi les plus robustes sur les
instances structurées (H8d résout les 17 instances SAT structurées).
Hors heuristique, 60 à 90\,\% du temps est passé dans la propagation, qui
recopie la formule à chaque simplification : la prochaine étape est une
structure de données sans recopie de la formule.

\section{Usage des assistants IA}

% TODO : à rédiger par le binôme (l'énoncé demande de le mentionner).

\end{document}
